\crefalias{section}{appendix}
\section{The Logic of Transition: Impedance Power Transfer}
\label{sec:impedancepowertransfer}

Unlike standard field theory, which treats interactions as point-like kinematic couplings, \textit{Informational Energetics} models them as \textbf{Impedance Power Transfer} events. To minimize the local dissipation density ($\mathcal{D} = \xi \cdot \mathbf{K} \cdot f$) defined in \cref{sec:dynamical_principle}, information flow across the lattice must strictly obey geometric impedance matching.

\subsection{The Admittance Amplitude (\texorpdfstring{$\sqrt{\alpha}$}{sqrtalpha})}
A topological defect (particle) maintains stability against the baseline vacuum impedance $Z_{vac} = \alpha^{-1}$

An interaction represents a \textbf{Topological Splice}, a geometric operation that conserves information density while allowing the knot to shed accumulated stress via a coordination signal (a boson). Because action (probability) scales with the inverse of impedance, the probability amplitude $\mathcal{A}$ of this transition event is defined strictly by the Admittance of the lattice node:
\begin{equation}
\mathcal{A} \propto \frac{1}{\sqrt{Z_{vac}}} = \sqrt{\alpha}
\end{equation}
This recovers the standard QED vertex factor ($e = \sqrt{4\pi\alpha}$), identifying electric charge not as an arbitrary parameter, but as the fundamental geometric admittance of the substrate.

\subsection{Dynamic Validation: Scattering Cross-Section}
To verify that this impedance logic holds dynamically, we derive the scattering cross-section for electron-positron annihilation ($e^+e^- \to \mu^+\mu^-$). In the lattice, this represents power transfer across the bulk interaction geometry.

\begin{enumerate}
    \item \textbf{Coupling Efficiency ($\eta_C$):} The total probability of the combined transition scales with the square of the input and output admittances: $\eta_C = (\sqrt{\alpha} \cdot \sqrt{\alpha})^2 = \alpha^2$.
    \item \textbf{Geometric Factor (1/3):} The signal propagates through the \textbf{Interaction Remainder} ($\sigma - \chi = 3$), representing the three spatial dimensions of the manifold bulk. Vector projection of the flux onto this 3D geometry introduces a strict normalization factor of $1/(\sigma - \chi) = 1/3$.
    \item \textbf{Aperture ($1/s$):} Unitarity requires the interaction area to scale with the inverse square of the center-of-mass energy ($1/s$).
\end{enumerate}

Combining these factors with the circular interface boundary ($4\pi$) yields:
\begin{equation}
\sigma_{E8} = \frac{4\pi \alpha^2}{3s}
\end{equation}
This exact recovery of the high-energy QED limit confirms that interaction probabilities are governed by geometric impedance matching. In the main text, we apply this exact same logic to the \textbf{Internal Transitions} (Flavor Mixing) between lattice resonance tiers.

\subsection{Geometric Validation: The Schwinger Limit}

Before applying this framework to the complexities of the CKM and PMNS matrices, we demonstrate its predictive power by recovering a foundational result of QED: the anomalous magnetic moment of the electron ($a_e$).

In standard QFT, this value arises from virtual particle loops in the vacuum state. The first-order correction, calculated by Julian Schwinger in 1948, is exactly:
\begin{equation}
a_e = \frac{\alpha}{2\pi} \approx 0.0011614
\end{equation}

Deriving this value in QFT requires evaluating the one-loop vertex correction diagram, a complex process involving renormalization, infrared regularization, and integration over Feynman parameters. In the $E_8$-Persistence theory, this is not a calculus problem; it is a \textbf{Geometric Aspect Ratio}.

\subsubsection{The Geometric Setup}

The electron is a geometric knot with a Unitary Address ($N=0$) coupled to the Gauge Connection (the photon). The interaction is defined by two invariant properties:

\begin{itemize}
    \item \textbf{The Flux ($\alpha$):} The strength of the interaction is governed by the Vacuum Admittance coefficient $\alpha$.
    \item \textbf{The Topology ($2\pi$):} While the electron spinor lives on the ``double cover" of the manifold ($4\pi$), the vector gauge field (photon) propagates strictly along the manifold boundary, which possesses a fundamental cycle of $2\pi$.
\end{itemize}

\subsubsection{The Derivation}

The anomalous moment represents the \textbf{Geometric Phase Lag} between the linear gauge field and the topological knot. It is simply the ratio of the Coupling Flux to the Gauge Cycle:

\begin{equation}
a_e = \frac{\text{Coupling Flux}}{\text{Gauge Cycle}} = \frac{\alpha}{2\pi}
\end{equation}

The geometric perspective explains why the phenomenologically complicated QFT calculus mathematically resolves to such a simple expression: the vacuum is merely enforcing an integer-bounded topological ratio.